% BCT Superfluid Lattice Model — Volume I: Foundations
% Letters L1–L12 | March 2026
% Michel Robert Cabrié | ORCID: 0009-0007-9561-9859
% Truncated Physics Press · Barrys Reef, Victoria, Australia
% Compile with: pdflatex → bibtex → pdflatex → pdflatex

\documentclass[aps,prl,twocolumn,showpacs,preprintnumbers]{revtex4-2}

\usepackage{amsmath,amssymb,amsfonts}
\usepackage{xcolor}
\usepackage{booktabs}
\usepackage{microtype}
\usepackage{hyperref}
\usepackage{geometry}

\definecolor{bctnavy}{RGB}{11,37,69}
\definecolor{bctblue}{RGB}{13,71,161}
\definecolor{bctteal}{RGB}{0,96,100}
\definecolor{bctgreen}{RGB}{27,94,32}

\hypersetup{
  colorlinks=true,
  linkcolor=bctnavy,
  citecolor=bctblue,
  urlcolor=bctteal
}

\newcommand{\roct}{r_{\mathrm{oct}}}
\newcommand{\rtet}{r_{\mathrm{tet}}}
\newcommand{\azero}{\alpha_0}
\newcommand{\LQCD}{\Lambda_{\mathrm{QCD}}}

\begin{document}

\preprint{BCT Letters Collection --- Volume I --- March 2026}

\title{BCT Letters Collection --- Volume I: Foundations\\
\large The BCT Vacuum Postulate, Geometry, and Electroweak Sector}

\author{Michel Robert Cabri\'e}
\affiliation{Independent Artist \& Researcher, Barrys Reef, Victoria, Australia (Pop.\ 28)}
\email{ZeroFreeParameters@gmail.com}
\homepage{https://patreon.com/cw/TheBCTSuperfluidLatticeModel}

\date{March 2026}

\begin{abstract}
Volume~I of the BCT Superfluid Lattice Model programme collects
the twelve foundational Letters establishing the geometric basis
of the entire framework. From three inputs --- $\roct = (\sqrt{2}-1)/2$,
$\rtet = (\sqrt{6}-2)/4$, and $\LQCD = 220$~MeV --- the BCT vacuum
geometry is derived, the D4 root lattice is identified as the unique
vacuum lattice consistent with Lorentz invariance, and the electromagnetic
and electroweak sectors are derived with zero free parameters.
Key results: fine structure constant $\alpha = \alpha_0(1-2\alpha_0)
= 1/137.018$ ($-0.013\%$); $\sin^2\theta_W = 0.23143$ ($+0.09\%$);
$m_e = 0.511$~MeV ($<0.01\%$). All results use zero free parameters.
Open access CC~BY~4.0.
\end{abstract}

\maketitle

%----------------------------------------------------------------------
\section*{Programme Context}
%----------------------------------------------------------------------

The BCT Superfluid Lattice Model proposes that the quantum vacuum is
a Planck-scale body-centred tetragonal superfluid crystal with axial
ratio $c/a = \sqrt{2}$, forced uniquely by Lorentz invariance. The
two void radii --- octahedral ($\roct$) and tetrahedral ($\rtet$) ---
are the only primitive length scales of physics. Every constant of
nature follows from their geometry.

This volume presents Letters~1--12, which establish the foundational
framework. The complete theoretical derivation is provided in the BCT
Monograph~\cite{monograph} and its companion appendix volumes~\cite{appvol1,appvol2}.
All letters in this volume are available individually on Zenodo under
ORCID~0009-0007-9561-9859.

\medskip
\noindent\textbf{Three inputs. Zero free parameters. Everything else: derived.}

\begin{equation}
\alpha_0 = \frac{\roct \cdot \rtet}{\pi}
= \frac{(\sqrt{2}-1)(\sqrt{6}-2)}{8\pi}
= 0.0074081
\end{equation}

%----------------------------------------------------------------------
\section{Letter 1 --- The BCT Vacuum Postulate}
%----------------------------------------------------------------------
\noindent\textbf{DOI:} \href{https://doi.org/10.5281/zenodo.18958207}{10.5281/zenodo.18958207}

\medskip
The quantum vacuum is a Planck-scale D4 superfluid crystal with
axial ratio $c/a = \sqrt{2}$. The two interstitial void radii
$\roct = (\sqrt{2}-1)/2$ and $\rtet = (\sqrt{6}-2)/4$ are the
only primitive length scales. Everything follows. This Letter
establishes the foundational postulate and derives the first
consequences: the BCT coupling $\azero$, the identification of
the D4 root lattice, and the emergence of three fermion generations
from D4 triality.

%----------------------------------------------------------------------
\section{Letter 2 --- The 24-Cell}
%----------------------------------------------------------------------
\noindent\textbf{DOI:} \href{https://doi.org/10.5281/zenodo.18958425}{10.5281/zenodo.18958425}

\medskip
The 24-cell polytope in four dimensions is identified as the
fundamental geometric object underlying the BCT vacuum. Its
24 vertices correspond to the 24 roots of the D4 root lattice.
The kissing number of D4 is 24, giving the instanton action
$S_{D4} = \pi^5/6$. The 24-cell is the unique self-dual regular
polytope in four dimensions --- it is its own dual, reflecting
the modular self-duality of the BCT vacuum.

%----------------------------------------------------------------------
\section{Letter 3 --- Octonions and the BCT Gauge Structure}
%----------------------------------------------------------------------
\noindent\textbf{DOI:} \href{https://doi.org/10.5281/zenodo.18958946}{10.5281/zenodo.18958946}

\medskip
The octonionic structure of the BCT gauge sector is derived. The
D4 root lattice admits a natural octonionic interpretation: the
seven imaginary units of the octonions correspond to the seven
non-trivial elements of the Fano plane, which is the projective
plane over $\mathbb{F}_2$. This structure forces the gauge group
to be $\mathrm{SU}(3) \times \mathrm{SU}(2) \times \mathrm{U}(1)$
--- the Standard Model gauge group --- as the unique group compatible
with D4 triality and octonionic multiplication.

%----------------------------------------------------------------------
\section{Letter 4 --- Injectivity}
%----------------------------------------------------------------------
\noindent\textbf{DOI:} \href{https://doi.org/10.5281/zenodo.18958769}{10.5281/zenodo.18958769}

\medskip
The injectivity of the BCT map from geometric inputs to Standard
Model observables is established. No two distinct geometric configurations
produce the same physical predictions. This result, combined with
the BCT Uniqueness Theorem (Letter~18), establishes that the BCT
geometry is the \emph{unique} solution: given Lorentz invariance
and the requirement of three fermion generations, the Standard Model
emerges necessarily.

%----------------------------------------------------------------------
\section{Letter 5 --- Vacuum Decay}
%----------------------------------------------------------------------
\noindent\textbf{DOI:} \href{https://doi.org/10.5281/zenodo.18960318}{10.5281/zenodo.18960318}

\medskip
The BCT vacuum is shown to be stable against decay. The bare
cosmological constant $\Lambda_{\mathrm{bare}} = 0$ exactly,
from the D4 modular self-duality. Vacuum decay requires tunnelling
through a potential barrier whose height diverges in the Planck
limit, making the BCT vacuum the unique stable ground state.
Prediction: the cosmological constant is not zero by fine-tuning
but by geometric necessity.

%----------------------------------------------------------------------
\section{Letter 6}
%----------------------------------------------------------------------
\noindent\textbf{DOI:} \href{https://doi.org/10.5281/zenodo.18959221}{10.5281/zenodo.18959221}

\medskip
The Maxwell equations are derived from BCT condensate geometry
using the Unruh--Visser acoustic analogy. The electromagnetic
field emerges as the long-wavelength limit of octahedral-void
phonon modes in the D4 condensate. The speed of light $c$ is
identified with the BCT phonon propagation speed, fixed uniquely
by the Lorentz invariance condition $c/a = \sqrt{2}$.

%----------------------------------------------------------------------
\section{Letter 7}
%----------------------------------------------------------------------
\noindent\textbf{DOI:} \href{https://doi.org/10.5281/zenodo.18959749}{10.5281/zenodo.18959749}

\medskip
The Dirac equation is derived from BCT vortex dynamics. Fermionic
excitations are identified with quantised vortices threading the
interstitial void network. The four-component Dirac spinor emerges
from the four-fold degeneracy of the D4 vortex core. The spin-statistics
theorem is a geometric consequence: vortices in a superfluid obey
Fermi statistics.

%----------------------------------------------------------------------
\section{Letter 8 --- The Complete CKM and PMNS Mixing Matrices}
%----------------------------------------------------------------------
\noindent\textbf{DOI:} \href{https://doi.org/10.5281/zenodo.18959548}{10.5281/zenodo.18959548}

\medskip
The CKM and PMNS mixing matrices are derived from the tetrahedral
void symmetry group $T_d$ (order 24). Key results:
\begin{itemize}
  \item Cabibbo angle: $\theta_C = \arctan(\rtet/\roct) = 12.87°$ (obs.\ $13.04°$, $-1.3\%$)
  \item CKM CP phase: $\delta_{\mathrm{CKM}} = \arccos(1/3) = 70.53°$ (tetrahedral vertex angle)
  \item PMNS CP phase: $\delta_{\mathrm{CP}} = -\pi/2$ exactly, from $T_d$ quaternion structure
  \item Reactor mixing: $\sin^2\theta_{13} = 3\alpha_0 = 0.0222$ ($+1.0\%$)
\end{itemize}
All from zero free parameters.

%----------------------------------------------------------------------
\section{Letter 9}
%----------------------------------------------------------------------
\noindent\textbf{DOI:} \href{https://doi.org/10.5281/zenodo.18959397}{10.5281/zenodo.18959397}

\medskip
Newton's constant is derived from the BCT condensate equation of
state. The gravitational constant $G = c^4/(8\pi K_{\mathrm{BCT}})$
where $K_{\mathrm{BCT}}$ is the bulk modulus of the BCT superfluid.
This identifies gravity as an emergent acoustic phenomenon --- not
a fundamental force but the long-range restoring pressure of the
Planck-scale condensate. The derivation reproduces $G$ to the
accuracy of the input Planck mass.

%----------------------------------------------------------------------
\section{Letter 10}
%----------------------------------------------------------------------
\noindent\textbf{DOI:} \href{https://doi.org/10.5281/zenodo.18958635}{10.5281/zenodo.18958635}

\medskip
Three fermion generations are derived from D4 triality. The outer
automorphism group of D4 is $S_3$ (the symmetric group on three
elements), which contains a $\mathbb{Z}_3$ subgroup --- the triality
automorphism. This $\mathbb{Z}_3$ permutes the three eight-dimensional
representations of D4: the vector $\mathbf{8}_v$, and the two spinors
$\mathbf{8}_s$ and $\mathbf{8}_c$. Exactly three generations of
fermions emerge, one per triality class. No fourth generation is
possible within D4 triality.

%----------------------------------------------------------------------
\section{Letter 11 --- Sixteen Structural Theorems for the Meson Spectrum}
%----------------------------------------------------------------------
\noindent\textbf{DOI:} \href{https://doi.org/10.5281/zenodo.18959915}{10.5281/zenodo.18959915}

\medskip
Sixteen parameter-free algebraic identities connect BCT void geometry
directly to meson masses, decay constants, and electromagnetic radii.
Key results:
\begin{itemize}
  \item $m_\rho = m_\pi/(2\sqrt{\alpha_0})(1-3\alpha_0/2) = 775.5$~MeV ($+0.035\%$)
  \item $m_K = 494.2$~MeV ($+0.1\%$)
  \item $m_\eta = 547.5$~MeV ($-0.074\%$)
  \item $m_\omega = 782.1$~MeV ($-0.070\%$)
  \item $r_p = 0.8386$~fm ($-0.33\%$), $r_\pi = 0.6580$~fm ($-0.16\%$)
\end{itemize}
All sixteen theorems falsifiable; each predicts a specific number
from geometry alone.

%----------------------------------------------------------------------
\section{Letter 12}
%----------------------------------------------------------------------
\noindent\textbf{DOI:} \href{https://doi.org/10.5281/zenodo.18960121}{10.5281/zenodo.18960121}

\medskip
The strong CP problem is resolved geometrically. Three independent
mechanisms enforce $\bar{\theta}_{\mathrm{QCD}} = 0$ exactly:
(i) the $\mathbb{Z}_2$ centre of the D4 Weyl group maps
$\theta \to -\theta$, requiring $\theta_{\mathrm{bare}} = 0$;
(ii) $T_d$ reflection symmetry forces all BCT Yukawa couplings
to be real; (iii) D4 modular self-duality makes the partition
function satisfy $Z(\theta) = Z(-\theta)$.
BCT predicts $d_n = 1.0 \times 10^{-32}$~$e$\,cm, testable at
next-generation neutron EDM experiments. No axion required or predicted.

%----------------------------------------------------------------------
\section*{Volume I Summary}
%----------------------------------------------------------------------

\begin{table}[h]
\centering
\caption{Selected key results from Volume~I}
\begin{tabular}{lcc}
\toprule
Observable & BCT & Error \\
\midrule
$1/\alpha$ & 137.018 & $-0.013\%$ \\
$\sin^2\theta_W$ & 0.23143 & $+0.09\%$ \\
$\delta_{\mathrm{CP}}$ (PMNS) & $-\pi/2$ & exact \\
$\sin^2\theta_{13}$ & 0.0222 & $+1.0\%$ \\
$m_\rho$ (MeV) & 775.5 & $+0.035\%$ \\
$m_K$ (MeV) & 494.2 & $+0.1\%$ \\
$\bar{\theta}_{\mathrm{QCD}}$ & 0 & exact \\
\bottomrule
\end{tabular}
\end{table}

\noindent All results from $\roct$, $\rtet$, $\LQCD$ only.
No fitting. No tuning. Zero free parameters.

%----------------------------------------------------------------------
\begin{thebibliography}{9}

\bibitem{monograph}
M.~R.\ Cabri\'e,
\textit{The BCT Superfluid Lattice Model: Complete Derivation of
Standard Model Parameters from Vacuum Geometry},
Zenodo (2026),
\href{https://doi.org/10.5281/zenodo.18884976}{10.5281/zenodo.18884976}.

\bibitem{appvol1}
M.~R.\ Cabri\'e,
\textit{Supplementary Material Volume 1: Technical Appendices ---
Foundations, Gauge, Mass, Precision},
Zenodo (2026),
\href{https://doi.org/10.5281/zenodo.18885133}{10.5281/zenodo.18885133}.

\bibitem{appvol2}
M.~R.\ Cabri\'e,
\textit{Supplementary Material Volume 2: Technical Appendices ---
Extended Precision, Hadronic Tower},
Zenodo (2026),
\href{https://doi.org/10.5281/zenodo.18885472}{10.5281/zenodo.18885472}.

\bibitem{uniqueness}
M.~R.\ Cabri\'e,
\textit{BCT Letter 18: The BCT Uniqueness Theorem},
Zenodo (2026),
\href{https://doi.org/10.5281/zenodo.18957202}{10.5281/zenodo.18957202}.

\bibitem{strongCP}
M.~R.\ Cabri\'e,
\textit{Geometric Resolution of the Strong CP Problem without an Axion},
Zenodo (2026),
\href{https://doi.org/10.5281/zenodo.18885628}{10.5281/zenodo.18885628}.

\end{thebibliography}

\end{document}
